\documentclass[11pt,a4paper]{article}

%=========================================================
% PACKAGES
%=========================================================
\usepackage[utf8]{inputenc}
\usepackage[T1]{fontenc}
\usepackage[french]{babel}
\usepackage[a4paper,margin=2.15cm]{geometry}
\usepackage{lmodern}
\usepackage{amsmath,amssymb,amsthm,mathtools}
\usepackage{fancyhdr}
\usepackage{lastpage}
\usepackage{enumitem}
\usepackage{array,tabularx}
\usepackage{tikz}
\usepackage[most]{tcolorbox}
\usepackage{hyperref}
\usepackage{xcolor}

%=========================================================
% MISE EN PAGE
%=========================================================
\setlength{\headheight}{14pt}
\pagestyle{fancy}
\fancyhf{}
\fancyhead[L]{Mathématiques CPGE}
\fancyhead[C]{Valeurs propres et vecteurs propres}
\fancyhead[R]{Page \thepage/\pageref{LastPage}}
\fancyfoot[C]{Cours complet -- Réduction des endomorphismes}

\hypersetup{
    colorlinks=true,
    linkcolor=blue!50!black,
    urlcolor=blue!50!black
}

\setlist[itemize]{label=\(\triangleright\)}
\setlist[enumerate]{label=\textbf{\arabic*.}}

%=========================================================
% COULEURS ET BOÎTES
%=========================================================
\definecolor{bleufonce}{RGB}{25,65,130}
\definecolor{bleuclair}{RGB}{235,243,255}
\definecolor{vertfonce}{RGB}{30,120,70}
\definecolor{vertclair}{RGB}{235,250,240}
\definecolor{rougefonce}{RGB}{170,40,40}
\definecolor{rougeclair}{RGB}{255,238,238}
\definecolor{orangefonce}{RGB}{190,110,20}
\definecolor{orangeclair}{RGB}{255,246,230}
\definecolor{violetfonce}{RGB}{95,55,145}
\definecolor{violetclair}{RGB}{246,241,255}

\tcbset{
    enhanced,
    sharp corners,
    boxrule=0.75pt,
    before skip=8pt,
    after skip=8pt
}

\newtcolorbox{definitionbox}[1][]{
    colback=bleuclair,
    colframe=bleufonce,
    title=\textbf{Définition #1}
}

\newtcolorbox{theorembox}[1][]{
    colback=vertclair,
    colframe=vertfonce,
    title=\textbf{Théorème / Propriété #1}
}

\newtcolorbox{methodbox}[1][]{
    colback=orangeclair,
    colframe=orangefonce,
    title=\textbf{Méthode #1}
}

\newtcolorbox{retenirbox}[1][]{
    colback=violetclair,
    colframe=violetfonce,
    title=\textbf{À retenir #1}
}

\newtcolorbox{erreurbox}[1][]{
    colback=rougeclair,
    colframe=rougefonce,
    title=\textbf{Erreur classique #1}
}

\newtcolorbox{exercicebox}[1][]{
    colback=white,
    colframe=black!60,
    title=\textbf{Exercice #1}
}

\newtcolorbox{correctionbox}[1][]{
    colback=gray!5,
    colframe=black!65,
    title=\textbf{Correction #1}
}

%=========================================================
% COMMANDES
%=========================================================
\newcommand{\K}{\mathbb K}
\newcommand{\R}{\mathbb R}
\newcommand{\C}{\mathbb C}
\newcommand{\N}{\mathbb N}
\newcommand{\M}{\mathcal M}
\newcommand{\Lcal}{\mathcal L}
\newcommand{\Vect}{\operatorname{Vect}}
\newcommand{\Ker}{\operatorname{Ker}}
\newcommand{\Imf}{\operatorname{Im}}
\newcommand{\rg}{\operatorname{rg}}
\newcommand{\tr}{\operatorname{tr}}
\newcommand{\Id}{\operatorname{Id}}
\newcommand{\Sp}{\operatorname{Sp}}
\newcommand{\diag}{\operatorname{diag}}
\newcommand{\Mat}{\operatorname{Mat}}
\newcommand{\Card}{\operatorname{Card}}

\newtheorem{prop}{Proposition}[section]
\newtheorem{lemme}[prop]{Lemme}
\newtheorem{corollaire}[prop]{Corollaire}

%=========================================================
% DOCUMENT
%=========================================================
\begin{document}

\begin{center}
    {\Huge \textbf{Valeurs propres et vecteurs propres}}\\[2mm]
    {\Large Directions invariantes, polynôme caractéristique et premiers critères de diagonalisation}\\[2mm]
    {\large Cours complet pour classes préparatoires scientifiques}\\[2mm]
    \rule{0.88\linewidth}{0.5pt}
\end{center}

\vspace{0.4cm}

\tableofcontents

\newpage

%=========================================================
\section*{Introduction générale}
\addcontentsline{toc}{section}{Introduction générale}

L'objectif de ce chapitre est de comprendre une idée centrale de l'algèbre linéaire :

\begin{center}
\textit{Existe-t-il des directions de l'espace qui sont conservées par un endomorphisme ?}
\end{center}

Soit \(E\) un espace vectoriel sur un corps \(\K\), et soit :
\[
u\in\Lcal(E).
\]
Un vecteur quelconque \(x\in E\) peut être envoyé par \(u\) dans une direction complètement différente.  
Mais certains vecteurs non nuls ont une propriété remarquable : leur image reste colinéaire à eux-mêmes.

Autrement dit, il peut exister :
\[
x\neq0
\]
et :
\[
\lambda\in\K
\]
tels que :
\[
u(x)=\lambda x.
\]

Dans ce cas, la droite vectorielle :
\[
\Vect(x)
\]
est stable par \(u\), et \(u\) agit sur cette droite comme une homothétie de rapport \(\lambda\).

\begin{retenirbox}
Un vecteur propre est un vecteur non nul dont la direction est conservée par l'endomorphisme.  
La valeur propre mesure le facteur de dilatation ou de contraction sur cette direction.
\end{retenirbox}

Ce chapitre est la porte d'entrée vers la réduction des endomorphismes :
\begin{itemize}
    \item calcul des valeurs propres ;
    \item calcul des sous-espaces propres ;
    \item diagonalisation ;
    \item calcul de puissances de matrices ;
    \item suites vectorielles ;
    \item polynômes annulateurs ;
    \item étude géométrique des transformations linéaires.
\end{itemize}

%=========================================================
\section{Pré-requis algébriques indispensables}

\subsection{Endomorphismes}

\begin{definitionbox}[Endomorphisme]
Soit \(E\) un espace vectoriel sur \(\K\).  
Un \textbf{endomorphisme} de \(E\) est une application linéaire :
\[
u:E\to E.
\]
On note :
\[
\Lcal(E)
\]
l'ensemble des endomorphismes de \(E\).
\end{definitionbox}

\begin{retenirbox}
Les valeurs propres concernent les endomorphismes, c'est-à-dire les applications linéaires dont l'espace de départ et l'espace d'arrivée sont les mêmes.
\end{retenirbox}

\subsection{Noyau et injectivité}

\begin{definitionbox}[Noyau]
Soit \(u:E\to F\) une application linéaire.  
Le noyau de \(u\) est :
\[
\Ker(u)=\{x\in E:\ u(x)=0\}.
\]
\end{definitionbox}

\begin{theorembox}
Une application linéaire \(u:E\to F\) est injective si et seulement si :
\[
\Ker(u)=\{0\}.
\]
\end{theorembox}

\begin{proof}
Supposons \(u\) injective.  
Si \(x\in\Ker(u)\), alors :
\[
u(x)=0=u(0).
\]
Par injectivité :
\[
x=0.
\]
Donc :
\[
\Ker(u)=\{0\}.
\]

Réciproquement, supposons :
\[
\Ker(u)=\{0\}.
\]
Si :
\[
u(x)=u(y),
\]
alors :
\[
u(x-y)=0.
\]
Donc :
\[
x-y\in\Ker(u).
\]
Ainsi :
\[
x-y=0,
\]
d'où :
\[
x=y.
\]
Donc \(u\) est injective.
\end{proof}

\subsection{Matrices et changement de base}

\begin{definitionbox}[Matrice d'un endomorphisme]
Soit \(E\) un espace vectoriel de dimension \(n\), muni d'une base :
\[
\mathcal B=(e_1,\dots,e_n).
\]
La matrice de \(u\in\Lcal(E)\) dans la base \(\mathcal B\), notée :
\[
\Mat_{\mathcal B}(u),
\]
est la matrice dont la \(j\)-ième colonne est la colonne des coordonnées de \(u(e_j)\) dans la base \(\mathcal B\).
\end{definitionbox}

\begin{theorembox}[Changement de base]
Soient \(\mathcal B\) et \(\mathcal B'\) deux bases de \(E\).  
Soit \(P\) la matrice de passage de \(\mathcal B'\) vers \(\mathcal B\).  
Si :
\[
A=\Mat_{\mathcal B}(u),
\qquad
A'=\Mat_{\mathcal B'}(u),
\]
alors :
\[
A'=P^{-1}AP.
\]
\end{theorembox}

\begin{definitionbox}[Matrices semblables]
Deux matrices \(A,B\in\M_n(\K)\) sont dites semblables s'il existe \(P\in GL_n(\K)\) tel que :
\[
B=P^{-1}AP.
\]
\end{definitionbox}

\begin{retenirbox}
Deux matrices semblables représentent le même endomorphisme dans deux bases différentes.  
Les valeurs propres doivent donc être invariantes par similitude.
\end{retenirbox}

%=========================================================
\section{Définitions fondamentales}

\subsection{Valeur propre et vecteur propre}

\begin{definitionbox}[Valeur propre et vecteur propre]
Soit \(E\) un \(\K\)-espace vectoriel, et soit :
\[
u\in\Lcal(E).
\]
Un scalaire \(\lambda\in\K\) est une \textbf{valeur propre} de \(u\) s'il existe un vecteur non nul \(x\in E\) tel que :
\[
u(x)=\lambda x.
\]
Un tel vecteur \(x\) est appelé un \textbf{vecteur propre} de \(u\) associé à la valeur propre \(\lambda\).
\end{definitionbox}

\begin{erreurbox}
Le vecteur nul n'est jamais un vecteur propre.  
En effet :
\[
u(0)=\lambda 0
\]
pour tout \(\lambda\), donc le vecteur nul ne donne aucune information sur l'endomorphisme.
\end{erreurbox}

\begin{retenirbox}
\[
u(x)=\lambda x
\]
signifie que \(u\) ne change pas la direction de \(x\).  
Il peut seulement :
\begin{itemize}
    \item l'étirer si \(|\lambda|>1\) ;
    \item le contracter si \(|\lambda|<1\) ;
    \item le conserver si \(\lambda=1\) ;
    \item le renverser si \(\lambda<0\) dans le cas réel ;
    \item l'annuler si \(\lambda=0\).
\end{itemize}
\end{retenirbox}

\subsection{Interprétation géométrique}

Si \(x\neq0\) est un vecteur propre associé à \(\lambda\), alors la droite :
\[
D=\Vect(x)
\]
est stable par \(u\).  
En effet, pour tout \(t\in\K\),
\[
u(tx)=t u(x)=t\lambda x=\lambda(tx)\in D.
\]

Ainsi, chercher les vecteurs propres revient à chercher les droites vectorielles stables par l'endomorphisme.

\begin{theorembox}
Un vecteur non nul \(x\) est vecteur propre de \(u\) si et seulement si la droite \(\Vect(x)\) est stable par \(u\).
\end{theorembox}

\begin{proof}
Supposons que \(x\neq0\) soit vecteur propre.  
Il existe \(\lambda\in\K\) tel que :
\[
u(x)=\lambda x.
\]
Alors, pour tout \(t\in\K\),
\[
u(tx)=t u(x)=t\lambda x\in\Vect(x).
\]
Donc \(\Vect(x)\) est stable.

Réciproquement, supposons que la droite \(\Vect(x)\), avec \(x\neq0\), soit stable par \(u\).  
Alors :
\[
u(x)\in\Vect(x).
\]
Il existe donc \(\lambda\in\K\) tel que :
\[
u(x)=\lambda x.
\]
Donc \(x\) est vecteur propre associé à \(\lambda\).
\end{proof}

\subsection{Premiers exemples}

\begin{exercicebox}[Exemple fondamental]
Soit :
\[
u:\R^2\to\R^2,
\qquad
u(x,y)=(2x,3y).
\]
Déterminer des valeurs propres et des vecteurs propres.
\end{exercicebox}

\begin{correctionbox}
On observe :
\[
u(1,0)=(2,0)=2(1,0).
\]
Donc \(2\) est valeur propre, et \((1,0)\) est un vecteur propre associé.

De même :
\[
u(0,1)=(0,3)=3(0,1).
\]
Donc \(3\) est valeur propre, et \((0,1)\) est un vecteur propre associé.

Plus généralement :
\[
u(x,0)=(2x,0)=2(x,0),
\]
donc tous les vecteurs non nuls de la droite \(\R(1,0)\) sont propres pour la valeur propre \(2\).

Et :
\[
u(0,y)=(0,3y)=3(0,y),
\]
donc tous les vecteurs non nuls de la droite \(\R(0,1)\) sont propres pour la valeur propre \(3\).
\end{correctionbox}

\begin{exercicebox}[Identité et homothétie]
Soit \(E\) un \(\K\)-espace vectoriel non nul.
\begin{enumerate}
    \item Déterminer les valeurs propres de \(\Id_E\).
    \item Déterminer les valeurs propres de \(\alpha\Id_E\), où \(\alpha\in\K\).
\end{enumerate}
\end{exercicebox}

\begin{correctionbox}
Pour l'identité :
\[
\Id_E(x)=x=1\cdot x
\]
pour tout \(x\neq0\).  
Donc tous les vecteurs non nuls sont propres, associés à la valeur propre \(1\).  
L'unique valeur propre est donc \(1\).

Pour l'homothétie :
\[
(\alpha\Id_E)(x)=\alpha x
\]
pour tout \(x\neq0\).  
Donc tous les vecteurs non nuls sont propres, associés à la valeur propre \(\alpha\).  
L'unique valeur propre est donc \(\alpha\).
\end{correctionbox}

%=========================================================
\section{Sous-espaces propres et spectre}

\subsection{Sous-espace propre}

\begin{definitionbox}[Sous-espace propre]
Soit \(u\in\Lcal(E)\) et soit \(\lambda\in\K\).  
On appelle \textbf{sous-espace propre associé à \(\lambda\)} l'ensemble :
\[
E_\lambda(u)=\Ker(u-\lambda\Id_E).
\]
Autrement dit :
\[
E_\lambda(u)=\{x\in E:\ u(x)=\lambda x\}.
\]
\end{definitionbox}

\begin{theorembox}
Pour tout \(\lambda\in\K\), \(E_\lambda(u)\) est un sous-espace vectoriel de \(E\).
\end{theorembox}

\begin{proof}
On a :
\[
E_\lambda(u)=\Ker(u-\lambda\Id_E).
\]
Or \(u-\lambda\Id_E\) est une application linéaire de \(E\) dans \(E\).  
Son noyau est donc un sous-espace vectoriel de \(E\).  
Ainsi :
\[
E_\lambda(u)
\]
est un sous-espace vectoriel.
\end{proof}

\begin{theorembox}
Un scalaire \(\lambda\in\K\) est une valeur propre de \(u\) si et seulement si :
\[
E_\lambda(u)\neq\{0\}.
\]
\end{theorembox}

\begin{proof}
Par définition, \(\lambda\) est une valeur propre s'il existe un vecteur non nul \(x\) tel que :
\[
u(x)=\lambda x.
\]
Cela équivaut à :
\[
(u-\lambda\Id_E)(x)=0,
\]
avec \(x\neq0\).  
Autrement dit :
\[
x\in E_\lambda(u)
\quad\text{et}\quad
x\neq0.
\]
Donc :
\[
E_\lambda(u)\neq\{0\}.
\]
La réciproque est immédiate.
\end{proof}

\begin{retenirbox}
Les vecteurs propres associés à \(\lambda\), auxquels on ajoute le vecteur nul, forment le sous-espace propre :
\[
E_\lambda(u).
\]
\end{retenirbox}

\begin{erreurbox}
Un sous-espace propre contient toujours \(0\), mais \(0\) n'est jamais un vecteur propre.  
On dit donc :
\[
E_\lambda(u)\setminus\{0\}
\]
est l'ensemble des vecteurs propres associés à \(\lambda\).
\end{erreurbox}

\subsection{Spectre}

\begin{definitionbox}[Spectre]
L'ensemble des valeurs propres de \(u\) est appelé le \textbf{spectre} de \(u\), et se note :
\[
\Sp(u).
\]
Donc :
\[
\Sp(u)=\{\lambda\in\K:\ E_\lambda(u)\neq\{0\}\}.
\]
\end{definitionbox}

Pour une matrice \(A\in\M_n(\K)\), on note de même :
\[
\Sp(A).
\]

\begin{erreurbox}
Le spectre dépend du corps de base.  
Une matrice réelle peut ne pas avoir de valeur propre réelle, mais avoir des valeurs propres complexes.
\end{erreurbox}

%=========================================================
\section{Traduction matricielle}

\subsection{Valeurs propres d'une matrice}

\begin{definitionbox}[Valeur propre d'une matrice]
Soit \(A\in\M_n(\K)\).  
Un scalaire \(\lambda\in\K\) est une valeur propre de \(A\) s'il existe une colonne non nulle \(X\in\K^n\) telle que :
\[
AX=\lambda X.
\]
Une telle colonne \(X\) est appelée vecteur propre de \(A\) associé à \(\lambda\).
\end{definitionbox}

L'équation :
\[
AX=\lambda X
\]
s'écrit :
\[
AX-\lambda X=0,
\]
donc :
\[
(A-\lambda I_n)X=0.
\]

\begin{theorembox}
Soit \(A\in\M_n(\K)\).  
Alors :
\[
\lambda\in\Sp(A)
\Longleftrightarrow
\Ker(A-\lambda I_n)\neq\{0\}.
\]
\end{theorembox}

\begin{proof}
C'est la traduction directe de :
\[
AX=\lambda X
\]
en :
\[
(A-\lambda I_n)X=0.
\]
L'existence d'une solution non nulle équivaut au fait que le noyau de \(A-\lambda I_n\) ne soit pas réduit à \(\{0\}\).
\end{proof}

\begin{methodbox}[Calculer un sous-espace propre]
Pour calculer \(E_\lambda(A)\) :
\begin{enumerate}
    \item écrire la matrice :
    \[
    A-\lambda I_n ;
    \]
    \item résoudre le système homogène :
    \[
    (A-\lambda I_n)X=0 ;
    \]
    \item écrire l'ensemble des solutions sous forme paramétrée ;
    \item extraire une base du noyau.
\end{enumerate}
\end{methodbox}

\subsection{Valeur propre et non-inversibilité}

\begin{theorembox}
Soit \(A\in\M_n(\K)\).  
Alors :
\[
\lambda\in\Sp(A)
\Longleftrightarrow
A-\lambda I_n\text{ n'est pas inversible}.
\]
\end{theorembox}

\begin{proof}
On sait :
\[
\lambda\in\Sp(A)
\Longleftrightarrow
\Ker(A-\lambda I_n)\neq\{0\}.
\]
Or, pour une matrice carrée, avoir un noyau non nul équivaut à ne pas être inversible.  
Donc :
\[
\lambda\in\Sp(A)
\Longleftrightarrow
A-\lambda I_n\text{ n'est pas inversible}.
\]
\end{proof}

\begin{corollaire}
Pour \(A\in\M_n(\K)\),
\[
\lambda\in\Sp(A)
\Longleftrightarrow
\det(A-\lambda I_n)=0.
\]
\end{corollaire}

\begin{proof}
Une matrice carrée est non inversible si et seulement si son déterminant est nul.
\end{proof}

\subsection{Invariance par similitude}

\begin{theorembox}
Deux matrices semblables ont les mêmes valeurs propres.
\end{theorembox}

\begin{proof}
Soient \(A,B\in\M_n(\K)\) deux matrices semblables.  
Il existe \(P\in GL_n(\K)\) tel que :
\[
B=P^{-1}AP.
\]
Alors :
\[
B-\lambda I_n
=
P^{-1}AP-\lambda P^{-1}P
=
P^{-1}(A-\lambda I_n)P.
\]
Comme \(P\) est inversible, la matrice \(B-\lambda I_n\) est inversible si et seulement si \(A-\lambda I_n\) est inversible.  
Donc :
\[
\lambda\in\Sp(B)
\Longleftrightarrow
\lambda\in\Sp(A).
\]
Ainsi :
\[
\Sp(A)=\Sp(B).
\]
\end{proof}

\begin{retenirbox}
Les valeurs propres ne dépendent pas de la base choisie.  
Elles appartiennent à l'endomorphisme, pas à une matrice particulière.
\end{retenirbox}

%=========================================================
\section{Polynôme caractéristique}

\subsection{Définition}

\begin{definitionbox}[Polynôme caractéristique d'une matrice]
Soit \(A\in\M_n(\K)\).  
Le polynôme caractéristique de \(A\) est :
\[
\chi_A(X)=\det(XI_n-A).
\]
\end{definitionbox}

\begin{retenirbox}
Dans ce cours, on adopte la convention :
\[
\chi_A(X)=\det(XI_n-A).
\]
Avec cette convention, \(\chi_A\) est un polynôme unitaire de degré \(n\).
\end{retenirbox}

\begin{theorembox}
Si \(A\in\M_n(\K)\), alors \(\chi_A\) est un polynôme unitaire de degré \(n\).
\end{theorembox}

\begin{proof}
La matrice \(XI_n-A\) a pour coefficients diagonaux :
\[
X-a_{11},\dots,X-a_{nn}.
\]
Dans le développement du déterminant, le seul terme pouvant produire \(X^n\) est le produit des termes diagonaux :
\[
(X-a_{11})\cdots(X-a_{nn}).
\]
Le coefficient de \(X^n\) est donc \(1\).  
Ainsi \(\chi_A\) est unitaire de degré \(n\).
\end{proof}

\subsection{Polynôme caractéristique d'un endomorphisme}

\begin{definitionbox}[Polynôme caractéristique d'un endomorphisme]
Soit \(E\) un \(\K\)-espace vectoriel de dimension finie \(n\), et soit :
\[
u\in\Lcal(E).
\]
On définit :
\[
\chi_u(X)=\chi_A(X),
\]
où \(A\) est la matrice de \(u\) dans une base quelconque de \(E\).
\end{definitionbox}

\begin{theorembox}
La définition de \(\chi_u\) ne dépend pas de la base choisie.
\end{theorembox}

\begin{proof}
Soient \(A\) et \(B\) les matrices de \(u\) dans deux bases différentes.  
Alors \(A\) et \(B\) sont semblables.  
Il existe donc \(P\in GL_n(\K)\) tel que :
\[
B=P^{-1}AP.
\]
Alors :
\[
XI_n-B
=
XI_n-P^{-1}AP
=
P^{-1}(XI_n-A)P.
\]
Donc :
\[
\det(XI_n-B)
=
\det(P^{-1})\det(XI_n-A)\det(P).
\]
Or :
\[
\det(P^{-1})\det(P)=1.
\]
Ainsi :
\[
\det(XI_n-B)=\det(XI_n-A).
\]
Donc :
\[
\chi_B=\chi_A.
\]
\end{proof}

\subsection{Valeurs propres et racines du polynôme caractéristique}

\begin{theorembox}[Caractérisation par le polynôme caractéristique]
Soit \(A\in\M_n(\K)\).  
Alors :
\[
\lambda\in\Sp(A)
\Longleftrightarrow
\chi_A(\lambda)=0.
\]
\end{theorembox}

\begin{proof}
On a :
\[
\lambda\in\Sp(A)
\Longleftrightarrow
A-\lambda I_n\text{ n'est pas inversible}.
\]
Cela équivaut à :
\[
\det(A-\lambda I_n)=0.
\]
Or :
\[
\chi_A(\lambda)=\det(\lambda I_n-A).
\]
Comme :
\[
\lambda I_n-A=-(A-\lambda I_n),
\]
on a :
\[
\det(\lambda I_n-A)=(-1)^n\det(A-\lambda I_n).
\]
Donc :
\[
\det(A-\lambda I_n)=0
\Longleftrightarrow
\chi_A(\lambda)=0.
\]
Ainsi :
\[
\lambda\in\Sp(A)
\Longleftrightarrow
\chi_A(\lambda)=0.
\]
\end{proof}

\begin{corollaire}
Les valeurs propres d'une matrice sont exactement les racines de son polynôme caractéristique dans le corps de base.
\end{corollaire}

\begin{corollaire}
Une matrice de taille \(n\) possède au plus \(n\) valeurs propres distinctes dans \(\K\).
\end{corollaire}

\begin{proof}
Le polynôme caractéristique est de degré \(n\).  
Un polynôme non nul de degré \(n\) possède au plus \(n\) racines distinctes.
\end{proof}

\begin{erreurbox}
Le polynôme caractéristique peut ne pas être scindé sur le corps de base.  
Par exemple :
\[
A=
\begin{pmatrix}
0&-1\\
1&0
\end{pmatrix}
\]
a pour polynôme caractéristique :
\[
\chi_A(X)=X^2+1.
\]
Sur \(\R\), il n'a aucune valeur propre.  
Sur \(\C\), ses valeurs propres sont \(i\) et \(-i\).
\end{erreurbox}

%=========================================================
\section{Premiers exemples complets}

\subsection{Matrice diagonale}

\begin{exercicebox}[Matrice diagonale]
Soit :
\[
A=
\begin{pmatrix}
2&0&0\\
0&3&0\\
0&0&3
\end{pmatrix}.
\]
Déterminer les valeurs propres et les sous-espaces propres de \(A\).
\end{exercicebox}

\begin{correctionbox}
La matrice est diagonale, donc :
\[
\chi_A(X)=(X-2)(X-3)^2.
\]
Les valeurs propres sont :
\[
2
\quad\text{et}\quad
3.
\]

Pour \(\lambda=2\) :
\[
A-2I_3=
\begin{pmatrix}
0&0&0\\
0&1&0\\
0&0&1
\end{pmatrix}.
\]
Le système donne :
\[
y=0,\qquad z=0.
\]
Donc :
\[
E_2(A)=\Vect
\begin{pmatrix}
1\\0\\0
\end{pmatrix}.
\]

Pour \(\lambda=3\) :
\[
A-3I_3=
\begin{pmatrix}
-1&0&0\\
0&0&0\\
0&0&0
\end{pmatrix}.
\]
Le système donne :
\[
x=0.
\]
Donc :
\[
E_3(A)=\Vect
\left(
\begin{pmatrix}
0\\1\\0
\end{pmatrix},
\begin{pmatrix}
0\\0\\1
\end{pmatrix}
\right).
\]
\end{correctionbox}

\subsection{Matrice triangulaire}

\begin{exercicebox}[Matrice triangulaire]
Soit :
\[
A=
\begin{pmatrix}
1&1\\
0&2
\end{pmatrix}.
\]
Déterminer les valeurs propres et les sous-espaces propres de \(A\).
\end{exercicebox}

\begin{correctionbox}
Comme \(A\) est triangulaire :
\[
\chi_A(X)=(X-1)(X-2).
\]
Les valeurs propres sont :
\[
1
\quad\text{et}\quad
2.
\]

Pour \(\lambda=1\) :
\[
A-I_2=
\begin{pmatrix}
0&1\\
0&1
\end{pmatrix}.
\]
Le système donne :
\[
y=0.
\]
Donc :
\[
E_1(A)=\Vect
\begin{pmatrix}
1\\0
\end{pmatrix}.
\]

Pour \(\lambda=2\) :
\[
A-2I_2=
\begin{pmatrix}
-1&1\\
0&0
\end{pmatrix}.
\]
Le système donne :
\[
-y+x=0,
\]
c'est-à-dire :
\[
y=x.
\]
Donc :
\[
E_2(A)=\Vect
\begin{pmatrix}
1\\1
\end{pmatrix}.
\]
\end{correctionbox}

\subsection{Matrice réelle sans valeur propre réelle}

\begin{exercicebox}[Rotation du plan]
Soit :
\[
A=
\begin{pmatrix}
0&-1\\
1&0
\end{pmatrix}.
\]
Étudier ses valeurs propres sur \(\R\), puis sur \(\C\).
\end{exercicebox}

\begin{correctionbox}
On calcule :
\[
\chi_A(X)=
\det
\begin{pmatrix}
X&1\\
-1&X
\end{pmatrix}
=
X^2+1.
\]
Sur \(\R\), le polynôme \(X^2+1\) n'a aucune racine.  
Donc \(A\) n'a aucune valeur propre réelle.

Sur \(\C\), on a :
\[
X^2+1=(X-i)(X+i).
\]
Les valeurs propres complexes sont donc :
\[
i
\quad\text{et}\quad
-i.
\]
\end{correctionbox}

%=========================================================
\section{Valeurs propres des matrices particulières}

\subsection{Matrices triangulaires}

\begin{theorembox}
Les valeurs propres d'une matrice triangulaire sont ses coefficients diagonaux.
\end{theorembox}

\begin{proof}
Soit \(A=(a_{ij})\) une matrice triangulaire supérieure.  
Alors \(XI_n-A\) est encore triangulaire supérieure, de coefficients diagonaux :
\[
X-a_{11},\dots,X-a_{nn}.
\]
Son déterminant vaut le produit des coefficients diagonaux :
\[
\chi_A(X)=\prod_{k=1}^n(X-a_{kk}).
\]
Les racines du polynôme caractéristique sont donc les coefficients diagonaux de \(A\).
\end{proof}

\begin{erreurbox}
Les coefficients diagonaux d'une matrice quelconque ne sont pas forcément ses valeurs propres.  
Ce résultat est vrai pour les matrices triangulaires, pas pour toutes les matrices.
\end{erreurbox}

\subsection{Matrices diagonales}

\begin{corollaire}
Les valeurs propres d'une matrice diagonale sont ses coefficients diagonaux.
\end{corollaire}

\begin{proof}
Une matrice diagonale est en particulier triangulaire.
\end{proof}

\subsection{Projecteurs}

\begin{definitionbox}[Projecteur]
Un endomorphisme \(p\in\Lcal(E)\) est un projecteur si :
\[
p^2=p.
\]
\end{definitionbox}

\begin{theorembox}
Si \(p\) est un projecteur, alors ses valeurs propres appartiennent à :
\[
\{0,1\}.
\]
\end{theorembox}

\begin{proof}
Soit \(\lambda\) une valeur propre de \(p\), et soit \(x\neq0\) un vecteur propre associé.  
Alors :
\[
p(x)=\lambda x.
\]
En appliquant \(p\) une deuxième fois :
\[
p^2(x)=p(\lambda x)=\lambda p(x)=\lambda^2x.
\]
Mais \(p^2=p\), donc :
\[
p^2(x)=p(x)=\lambda x.
\]
Ainsi :
\[
\lambda^2x=\lambda x.
\]
Donc :
\[
\lambda(\lambda-1)x=0.
\]
Comme \(x\neq0\), on obtient :
\[
\lambda(\lambda-1)=0.
\]
Donc :
\[
\lambda\in\{0,1\}.
\]
\end{proof}

\begin{theorembox}
Pour un projecteur \(p\),
\[
E_1(p)=\Imf(p),
\qquad
E_0(p)=\Ker(p).
\]
\end{theorembox}

\begin{proof}
On a :
\[
E_0(p)=\Ker(p-0\Id)=\Ker(p).
\]

Montrons que :
\[
E_1(p)=\Imf(p).
\]
Si \(x\in E_1(p)\), alors :
\[
p(x)=x.
\]
Donc \(x\in\Imf(p)\), car \(x=p(x)\).

Réciproquement, si \(x\in\Imf(p)\), il existe \(y\in E\) tel que :
\[
x=p(y).
\]
Alors :
\[
p(x)=p(p(y))=p^2(y)=p(y)=x.
\]
Donc :
\[
x\in E_1(p).
\]
\end{proof}

\subsection{Involutions}

\begin{definitionbox}[Involution]
Un endomorphisme \(s\in\Lcal(E)\) est une involution si :
\[
s^2=\Id_E.
\]
\end{definitionbox}

\begin{theorembox}
Si \(s^2=\Id_E\), alors ses valeurs propres appartiennent à :
\[
\{-1,1\}.
\]
\end{theorembox}

\begin{proof}
Soit \(x\neq0\) un vecteur propre associé à \(\lambda\).  
Alors :
\[
s(x)=\lambda x.
\]
Donc :
\[
s^2(x)=\lambda^2x.
\]
Mais :
\[
s^2=\Id_E,
\]
donc :
\[
s^2(x)=x.
\]
Ainsi :
\[
\lambda^2x=x.
\]
Donc :
\[
(\lambda^2-1)x=0.
\]
Comme \(x\neq0\), on obtient :
\[
\lambda^2-1=0.
\]
Donc :
\[
\lambda\in\{-1,1\}.
\]
\end{proof}

\subsection{Matrices nilpotentes}

\begin{definitionbox}[Matrice nilpotente]
Une matrice \(A\in\M_n(\K)\) est nilpotente s'il existe \(k\geq1\) tel que :
\[
A^k=0.
\]
\end{definitionbox}

\begin{theorembox}
Si \(A\) est nilpotente, alors sa seule valeur propre possible est :
\[
0.
\]
\end{theorembox}

\begin{proof}
Soit \(\lambda\) une valeur propre de \(A\), et soit \(X\neq0\) un vecteur propre associé :
\[
AX=\lambda X.
\]
Alors, par récurrence :
\[
A^kX=\lambda^kX.
\]
Si \(A^k=0\), alors :
\[
0=A^kX=\lambda^kX.
\]
Comme \(X\neq0\), on obtient :
\[
\lambda^k=0.
\]
Donc :
\[
\lambda=0.
\]
\end{proof}

%=========================================================
\section{Indépendance des vecteurs propres}

\subsection{Théorème fondamental}

\begin{theorembox}[Liberté des vecteurs propres]
Des vecteurs propres associés à des valeurs propres deux à deux distinctes sont linéairement indépendants.
\end{theorembox}

\begin{proof}
Soient \(\lambda_1,\dots,\lambda_p\) des valeurs propres deux à deux distinctes de \(u\).  
Soient \(x_1,\dots,x_p\) des vecteurs propres associés respectivement à ces valeurs propres :
\[
u(x_i)=\lambda_i x_i,
\qquad
x_i\neq0.
\]

On montre que la famille \((x_1,\dots,x_p)\) est libre par récurrence sur \(p\).

Pour \(p=1\), la famille \((x_1)\) est libre car \(x_1\neq0\).

Supposons le résultat vrai au rang \(p-1\).  
Considérons une relation linéaire :
\[
\alpha_1x_1+\cdots+\alpha_px_p=0.
\]
En appliquant \(u\), on obtient :
\[
\alpha_1u(x_1)+\cdots+\alpha_pu(x_p)=0,
\]
donc :
\[
\alpha_1\lambda_1x_1+\cdots+\alpha_p\lambda_px_p=0.
\]
Multiplions la relation initiale par \(\lambda_p\) :
\[
\alpha_1\lambda_px_1+\cdots+\alpha_p\lambda_px_p=0.
\]
Soustrayons :
\[
\alpha_1(\lambda_1-\lambda_p)x_1
+\cdots+
\alpha_{p-1}(\lambda_{p-1}-\lambda_p)x_{p-1}=0.
\]
Par hypothèse de récurrence, la famille \((x_1,\dots,x_{p-1})\) est libre.  
Donc, pour tout \(i\in\{1,\dots,p-1\}\),
\[
\alpha_i(\lambda_i-\lambda_p)=0.
\]
Comme les valeurs propres sont deux à deux distinctes :
\[
\lambda_i-\lambda_p\neq0.
\]
Donc :
\[
\alpha_i=0
\]
pour tout \(i\leq p-1\).

La relation initiale devient :
\[
\alpha_px_p=0.
\]
Comme \(x_p\neq0\), on obtient :
\[
\alpha_p=0.
\]
Tous les coefficients sont nuls.  
La famille est donc libre.
\end{proof}

\begin{corollaire}
Les sous-espaces propres associés à des valeurs propres deux à deux distinctes sont en somme directe.
\end{corollaire}

\begin{proof}
Soient \(\lambda_1,\dots,\lambda_r\) des valeurs propres distinctes.  
Montrons que :
\[
E_{\lambda_1}\oplus\cdots\oplus E_{\lambda_r}
\]
est une somme directe.

Supposons :
\[
x_1+\cdots+x_r=0,
\]
avec :
\[
x_i\in E_{\lambda_i}.
\]
Si certains \(x_i\) sont non nuls, ils forment une famille de vecteurs propres associés à des valeurs propres distinctes.  
D'après le théorème précédent, cette famille est libre.  
La relation ci-dessus impose donc que tous les \(x_i\) soient nuls.

Ainsi la somme est directe.
\end{proof}

\begin{corollaire}
Si \(E\) est de dimension \(n\), alors \(u\) possède au plus \(n\) valeurs propres distinctes.
\end{corollaire}

\begin{proof}
Pour chaque valeur propre, choisissons un vecteur propre non nul.  
Si \(u\) avait \(p\) valeurs propres distinctes, on obtiendrait \(p\) vecteurs libres dans un espace de dimension \(n\).  
Donc :
\[
p\leq n.
\]
\end{proof}

\begin{retenirbox}
Des valeurs propres différentes produisent des directions indépendantes.  
C'est l'idée fondamentale qui rend la diagonalisation possible.
\end{retenirbox}

%=========================================================
\section{Multiplicités algébrique et géométrique}

\subsection{Multiplicité algébrique}

\begin{definitionbox}[Multiplicité algébrique]
Soit \(\lambda\in\Sp(u)\).  
La multiplicité algébrique de \(\lambda\), notée \(m_\lambda\), est sa multiplicité comme racine du polynôme caractéristique :
\[
\chi_u(X).
\]
\end{definitionbox}

\subsection{Multiplicité géométrique}

\begin{definitionbox}[Multiplicité géométrique]
La multiplicité géométrique de \(\lambda\) est :
\[
\dim(E_\lambda(u)).
\]
\end{definitionbox}

\begin{retenirbox}
La multiplicité algébrique se lit dans le polynôme caractéristique.  
La multiplicité géométrique se calcule en résolvant un système linéaire :
\[
E_\lambda=\Ker(u-\lambda\Id_E).
\]
\end{retenirbox}

\subsection{Inégalité fondamentale}

\begin{theorembox}[Inégalité fondamentale]
Pour toute valeur propre \(\lambda\),
\[
1\leq \dim(E_\lambda)\leq m_\lambda.
\]
\end{theorembox}

\begin{proof}
Comme \(\lambda\) est une valeur propre, on a :
\[
E_\lambda\neq\{0\}.
\]
Donc :
\[
\dim(E_\lambda)\geq1.
\]

Posons :
\[
r=\dim(E_\lambda).
\]
Choisissons une base :
\[
(e_1,\dots,e_r)
\]
de \(E_\lambda\), puis complétons-la en une base de \(E\) :
\[
(e_1,\dots,e_r,e_{r+1},\dots,e_n).
\]

Pour \(1\leq i\leq r\), on a :
\[
u(e_i)=\lambda e_i.
\]
Dans cette base, la matrice de \(u\) est donc de la forme :
\[
A=
\begin{pmatrix}
\lambda I_r & B\\
0 & C
\end{pmatrix}.
\]
Alors :
\[
XI_n-A=
\begin{pmatrix}
(X-\lambda)I_r & -B\\
0 & XI_{n-r}-C
\end{pmatrix}.
\]
En prenant le déterminant d'une matrice triangulaire par blocs :
\[
\chi_u(X)
=
\det(XI_n-A)
=
(X-\lambda)^r\det(XI_{n-r}-C).
\]
Donc \((X-\lambda)^r\) divise \(\chi_u(X)\).  
Cela signifie que :
\[
r\leq m_\lambda.
\]
Ainsi :
\[
1\leq \dim(E_\lambda)\leq m_\lambda.
\]
\end{proof}

\begin{erreurbox}
Il est impossible d'avoir :
\[
\dim(E_\lambda)>m_\lambda.
\]
Si cela arrive dans un calcul, il y a forcément une erreur dans le polynôme caractéristique ou dans le noyau.
\end{erreurbox}

%=========================================================
\section{Premiers critères de diagonalisation}

\subsection{Définition}

\begin{definitionbox}[Endomorphisme diagonalisable]
Un endomorphisme \(u\in\Lcal(E)\) est diagonalisable s'il existe une base de \(E\) formée de vecteurs propres de \(u\).
\end{definitionbox}

\begin{definitionbox}[Matrice diagonalisable]
Une matrice \(A\in\M_n(\K)\) est diagonalisable s'il existe une matrice inversible \(P\in GL_n(\K)\) et une matrice diagonale \(D\) telles que :
\[
A=PDP^{-1}.
\]
\end{definitionbox}

\begin{retenirbox}
Dans une diagonalisation :
\[
A=PDP^{-1},
\]
les colonnes de \(P\) sont des vecteurs propres de \(A\), et les coefficients diagonaux de \(D\) sont les valeurs propres correspondantes.
\end{retenirbox}

\subsection{Critère par somme directe}

\begin{theorembox}[Critère géométrique]
Un endomorphisme \(u\in\Lcal(E)\) est diagonalisable si et seulement si :
\[
E=\bigoplus_{\lambda\in\Sp(u)}E_\lambda(u).
\]
\end{theorembox}

\begin{proof}
Supposons \(u\) diagonalisable.  
Il existe une base de \(E\) formée de vecteurs propres de \(u\).  
En regroupant les vecteurs de cette base selon leur valeur propre, on obtient une base de chaque sous-espace propre.  
La réunion de ces bases forme une base de \(E\).  
Donc :
\[
E=\bigoplus_{\lambda\in\Sp(u)}E_\lambda(u).
\]

Réciproquement, supposons :
\[
E=\bigoplus_{\lambda\in\Sp(u)}E_\lambda(u).
\]
Pour chaque \(\lambda\), choisissons une base de \(E_\lambda(u)\).  
Comme la somme est directe et vaut \(E\), la réunion de ces bases est une base de \(E\).  
Tous les vecteurs de cette base sont des vecteurs propres.  
Donc \(u\) est diagonalisable.
\end{proof}

\begin{corollaire}[Critère dimensionnel]
Un endomorphisme \(u\) est diagonalisable si et seulement si :
\[
\sum_{\lambda\in\Sp(u)}\dim(E_\lambda(u))=\dim(E).
\]
\end{corollaire}

\begin{proof}
Les sous-espaces propres associés à des valeurs propres distinctes sont en somme directe.  
Donc la dimension de leur somme est :
\[
\sum_{\lambda\in\Sp(u)}\dim(E_\lambda(u)).
\]
Cette somme est égale à \(E\) si et seulement si sa dimension vaut \(\dim(E)\).
\end{proof}

\subsection{Critère des valeurs propres distinctes}

\begin{theorembox}
Si \(E\) est de dimension \(n\) et si \(u\) admet \(n\) valeurs propres distinctes dans \(\K\), alors \(u\) est diagonalisable.
\end{theorembox}

\begin{proof}
Choisissons un vecteur propre non nul associé à chacune des \(n\) valeurs propres.  
Ces \(n\) vecteurs propres sont associés à des valeurs propres deux à deux distinctes, donc ils forment une famille libre.

Une famille libre de \(n\) vecteurs dans un espace de dimension \(n\) est une base.  
On a donc une base de \(E\) formée de vecteurs propres.  
Ainsi \(u\) est diagonalisable.
\end{proof}

\subsection{Critère par les multiplicités}

\begin{theorembox}[Critère pratique de diagonalisation]
Un endomorphisme \(u\in\Lcal(E)\), avec \(E\) de dimension \(n\), est diagonalisable sur \(\K\) si et seulement si :
\begin{enumerate}
    \item le polynôme caractéristique \(\chi_u\) est scindé sur \(\K\) ;
    \item pour toute valeur propre \(\lambda\),
    \[
    \dim(E_\lambda)=m_\lambda.
    \]
\end{enumerate}
\end{theorembox}

\begin{proof}
Supposons \(u\) diagonalisable.  
Dans une base de vecteurs propres, sa matrice est diagonale :
\[
D=\diag(\lambda_1,\dots,\lambda_n).
\]
Alors :
\[
\chi_u(X)=\prod_{k=1}^n(X-\lambda_k),
\]
donc \(\chi_u\) est scindé.  
De plus, la multiplicité algébrique d'une valeur propre \(\lambda\) est exactement le nombre de fois où \(\lambda\) apparaît sur la diagonale, ce qui est aussi la dimension du sous-espace propre associé.  
Donc :
\[
\dim(E_\lambda)=m_\lambda.
\]

Réciproquement, supposons que \(\chi_u\) soit scindé et que :
\[
\dim(E_\lambda)=m_\lambda
\]
pour toute valeur propre \(\lambda\).  
Comme \(\chi_u\) est scindé de degré \(n\), la somme des multiplicités algébriques vaut :
\[
\sum_{\lambda\in\Sp(u)}m_\lambda=n.
\]
Donc :
\[
\sum_{\lambda\in\Sp(u)}\dim(E_\lambda)
=
\sum_{\lambda\in\Sp(u)}m_\lambda
=
n.
\]
Par le critère dimensionnel, \(u\) est diagonalisable.
\end{proof}

\begin{erreurbox}
Un polynôme caractéristique scindé ne suffit pas à garantir la diagonalisation.  
Il faut aussi vérifier les dimensions des sous-espaces propres.
\end{erreurbox}

%=========================================================
\section{Polynômes annulateurs et valeurs propres}

\subsection{Polynôme d'un endomorphisme}

\begin{definitionbox}[Polynôme d'un endomorphisme]
Soit :
\[
P(X)=a_0+a_1X+\cdots+a_mX^m\in\K[X].
\]
Pour \(u\in\Lcal(E)\), on définit :
\[
P(u)=a_0\Id_E+a_1u+\cdots+a_mu^m.
\]
\end{definitionbox}

\begin{definitionbox}[Polynôme annulateur]
Un polynôme \(P\in\K[X]\) est un polynôme annulateur de \(u\) si :
\[
P(u)=0.
\]
\end{definitionbox}

\subsection{Lien avec les valeurs propres}

\begin{theorembox}
Si \(P(u)=0\), alors toute valeur propre \(\lambda\) de \(u\) vérifie :
\[
P(\lambda)=0.
\]
\end{theorembox}

\begin{proof}
Soit \(x\neq0\) un vecteur propre associé à \(\lambda\).  
Alors :
\[
u(x)=\lambda x.
\]
Par récurrence :
\[
u^k(x)=\lambda^kx.
\]
Si :
\[
P(X)=a_0+a_1X+\cdots+a_mX^m,
\]
alors :
\[
P(u)(x)
=
a_0x+a_1u(x)+\cdots+a_mu^m(x).
\]
Donc :
\[
P(u)(x)
=
(a_0+a_1\lambda+\cdots+a_m\lambda^m)x
=
P(\lambda)x.
\]
Comme \(P(u)=0\), on a :
\[
0=P(u)(x)=P(\lambda)x.
\]
Comme \(x\neq0\), on obtient :
\[
P(\lambda)=0.
\]
\end{proof}

\begin{retenirbox}
Les valeurs propres de \(u\) sont toujours parmi les racines de tout polynôme annulateur de \(u\).
\[
P(u)=0
\quad\Longrightarrow\quad
\Sp(u)\subset\{\lambda\in\K:\ P(\lambda)=0\}.
\]
\end{retenirbox}

\begin{erreurbox}
Les racines d'un polynôme annulateur ne sont pas forcément toutes des valeurs propres.  
On a seulement une inclusion :
\[
\Sp(u)\subset\{\text{racines de }P\}.
\]
\end{erreurbox}

\subsection{Critère de diagonalisation par polynôme annulateur}

\begin{theorembox}
Si \(u\) admet un polynôme annulateur scindé à racines simples, alors \(u\) est diagonalisable.
\end{theorembox}

\begin{proof}
Supposons qu'il existe :
\[
P(X)=\prod_{i=1}^r(X-\lambda_i),
\]
avec les \(\lambda_i\) deux à deux distincts, tel que :
\[
P(u)=0.
\]
Les polynômes :
\[
X-\lambda_1,\dots,X-\lambda_r
\]
sont deux à deux premiers entre eux.

Par le théorème de décomposition des noyaux, on obtient :
\[
E=\Ker(P(u))
=
\bigoplus_{i=1}^r \Ker(u-\lambda_i\Id_E).
\]
Comme \(P(u)=0\), on a :
\[
\Ker(P(u))=E.
\]
Donc :
\[
E=\bigoplus_{i=1}^r E_{\lambda_i}(u).
\]
Ainsi \(E\) est somme directe de sous-espaces propres de \(u\).  
Donc \(u\) est diagonalisable.
\end{proof}

\begin{methodbox}[Utiliser un polynôme annulateur]
Pour exploiter une relation du type :
\[
P(u)=0,
\]
on suit le plan :
\begin{enumerate}
    \item factoriser \(P\) sur le corps de base ;
    \item en déduire les valeurs propres possibles ;
    \item si \(P\) est scindé à racines simples, conclure que \(u\) est diagonalisable ;
    \item sinon, calculer les sous-espaces propres pour conclure.
\end{enumerate}
\end{methodbox}

%=========================================================
\section{Méthode complète de recherche des éléments propres}

\begin{methodbox}[Méthode universelle]
Pour déterminer les valeurs propres et les sous-espaces propres d'une matrice \(A\in\M_n(\K)\) :
\begin{enumerate}
    \item Calculer :
    \[
    \chi_A(X)=\det(XI_n-A).
    \]
    \item Factoriser \(\chi_A\) sur le corps de base \(\K\).
    \item Résoudre :
    \[
    \chi_A(\lambda)=0.
    \]
    Les solutions sont les valeurs propres.
    \item Pour chaque valeur propre \(\lambda\), calculer :
    \[
    E_\lambda(A)=\Ker(A-\lambda I_n).
    \]
    \item Donner une base de chaque sous-espace propre.
    \item Comparer les dimensions des sous-espaces propres aux multiplicités algébriques si l'on veut étudier la diagonalisabilité.
\end{enumerate}
\end{methodbox}

\begin{methodbox}[Méthode rapide]
Si \(A\in\M_n(\K)\) possède \(n\) valeurs propres distinctes, alors elle est immédiatement diagonalisable.  
Il faut tout de même calculer les vecteurs propres si l'on demande une diagonalisation explicite.
\end{methodbox}

\begin{methodbox}[Méthode avec une matrice triangulaire]
Si \(A\) est triangulaire :
\begin{enumerate}
    \item les valeurs propres sont directement les coefficients diagonaux ;
    \item le polynôme caractéristique est :
    \[
    \chi_A(X)=\prod_{i=1}^n(X-a_{ii});
    \]
    \item il reste à calculer les noyaux :
    \[
    \Ker(A-\lambda I_n).
    \]
\end{enumerate}
\end{methodbox}

%=========================================================
\section{Exemples détaillés de difficulté croissante}

\subsection{Exemple 1 : matrice diagonalisable avec valeurs propres distinctes}

Soit :
\[
A=
\begin{pmatrix}
2&1\\
0&3
\end{pmatrix}.
\]

On calcule :
\[
\chi_A(X)=
\det
\begin{pmatrix}
X-2&-1\\
0&X-3
\end{pmatrix}
=
(X-2)(X-3).
\]
Les valeurs propres sont :
\[
2
\quad\text{et}\quad
3.
\]
Comme elles sont distinctes et que \(A\in\M_2(\R)\), \(A\) est diagonalisable.

Calculons les sous-espaces propres.

Pour \(\lambda=2\) :
\[
A-2I_2=
\begin{pmatrix}
0&1\\
0&1
\end{pmatrix}.
\]
Le système donne :
\[
y=0.
\]
Donc :
\[
E_2=\Vect
\begin{pmatrix}
1\\0
\end{pmatrix}.
\]

Pour \(\lambda=3\) :
\[
A-3I_2=
\begin{pmatrix}
-1&1\\
0&0
\end{pmatrix}.
\]
Le système donne :
\[
-x+y=0,
\]
donc :
\[
y=x.
\]
Ainsi :
\[
E_3=\Vect
\begin{pmatrix}
1\\1
\end{pmatrix}.
\]

On peut prendre :
\[
P=
\begin{pmatrix}
1&1\\
0&1
\end{pmatrix},
\qquad
D=
\begin{pmatrix}
2&0\\
0&3
\end{pmatrix}.
\]
Alors :
\[
A=PDP^{-1}.
\]

\subsection{Exemple 2 : polynôme scindé mais matrice non diagonalisable}

Soit :
\[
A=
\begin{pmatrix}
1&1\\
0&1
\end{pmatrix}.
\]
On a :
\[
\chi_A(X)=(X-1)^2.
\]
L'unique valeur propre est :
\[
1.
\]
Sa multiplicité algébrique vaut :
\[
m_1=2.
\]
Calculons :
\[
A-I_2=
\begin{pmatrix}
0&1\\
0&0
\end{pmatrix}.
\]
Le système :
\[
(A-I_2)X=0
\]
donne :
\[
y=0.
\]
Donc :
\[
E_1=\Vect
\begin{pmatrix}
1\\0
\end{pmatrix}.
\]
Ainsi :
\[
\dim(E_1)=1<2=m_1.
\]
La matrice n'est pas diagonalisable.

\subsection{Exemple 3 : matrice réelle sans valeurs propres réelles}

Soit :
\[
A=
\begin{pmatrix}
0&-1\\
1&0
\end{pmatrix}.
\]
On calcule :
\[
\chi_A(X)=X^2+1.
\]
Sur \(\R\), il n'y a pas de racine.  
Donc :
\[
\Sp_\R(A)=\varnothing.
\]
La matrice n'est pas diagonalisable sur \(\R\).

Sur \(\C\), les valeurs propres sont :
\[
i
\quad\text{et}\quad
-i.
\]
Elles sont distinctes, donc \(A\) est diagonalisable sur \(\C\).

\subsection{Exemple 4 : matrice \(3\times3\)}

Soit :
\[
A=
\begin{pmatrix}
2&1&0\\
0&2&0\\
0&0&3
\end{pmatrix}.
\]
La matrice est triangulaire, donc :
\[
\chi_A(X)=(X-2)^2(X-3).
\]
Les valeurs propres sont :
\[
2
\quad\text{et}\quad
3.
\]

Pour \(\lambda=2\) :
\[
A-2I_3=
\begin{pmatrix}
0&1&0\\
0&0&0\\
0&0&1
\end{pmatrix}.
\]
Le système donne :
\[
y=0,
\qquad
z=0.
\]
Donc :
\[
E_2=\Vect
\begin{pmatrix}
1\\0\\0
\end{pmatrix}.
\]
Ainsi :
\[
\dim(E_2)=1<2.
\]
La matrice n'est pas diagonalisable.

Pour \(\lambda=3\) :
\[
A-3I_3=
\begin{pmatrix}
-1&1&0\\
0&-1&0\\
0&0&0
\end{pmatrix}.
\]
Le système donne :
\[
x=0,\qquad y=0.
\]
Donc :
\[
E_3=\Vect
\begin{pmatrix}
0\\0\\1
\end{pmatrix}.
\]

La somme des dimensions vaut :
\[
1+1=2<3.
\]
Donc \(A\) n'est pas diagonalisable.

%=========================================================
\section{Applications simples}

\subsection{Calcul des puissances}

Si \(A\) est diagonalisable :
\[
A=PDP^{-1},
\]
alors :
\[
A^n=PD^nP^{-1}.
\]

\begin{exercicebox}[Puissances d'une matrice]
Soit :
\[
A=
\begin{pmatrix}
1&1\\
1&1
\end{pmatrix}.
\]
Calculer \(A^n\) pour \(n\geq1\).
\end{exercicebox}

\begin{correctionbox}
On calcule :
\[
\chi_A(X)=
\det
\begin{pmatrix}
X-1&-1\\
-1&X-1
\end{pmatrix}
=
(X-1)^2-1
=
X^2-2X
=
X(X-2).
\]
Les valeurs propres sont :
\[
0
\quad\text{et}\quad
2.
\]
Elles sont distinctes, donc \(A\) est diagonalisable.

Pour \(\lambda=0\) :
\[
E_0=\Vect
\begin{pmatrix}
1\\-1
\end{pmatrix}.
\]
Pour \(\lambda=2\) :
\[
E_2=\Vect
\begin{pmatrix}
1\\1
\end{pmatrix}.
\]

On prend :
\[
P=
\begin{pmatrix}
1&1\\
-1&1
\end{pmatrix},
\qquad
D=
\begin{pmatrix}
0&0\\
0&2
\end{pmatrix}.
\]
Alors :
\[
A=PDP^{-1}.
\]
Pour \(n\geq1\),
\[
D^n=
\begin{pmatrix}
0&0\\
0&2^n
\end{pmatrix}.
\]
Comme :
\[
P^{-1}
=
\frac12
\begin{pmatrix}
1&-1\\
1&1
\end{pmatrix},
\]
on obtient :
\[
A^n
=
PD^nP^{-1}
=
2^{n-1}
\begin{pmatrix}
1&1\\
1&1
\end{pmatrix}.
\]
\end{correctionbox}

\subsection{Suites vectorielles}

Si :
\[
X_{n+1}=AX_n,
\]
alors :
\[
X_n=A^nX_0.
\]
La diagonalisation permet de calculer \(A^n\), donc d'obtenir une expression explicite de la suite.

%=========================================================
\newpage
\section{Exercices progressifs}

\subsection*{Thème 1 -- Définitions et premières recherches}

\begin{exercicebox}[1]
Soit :
\[
A=
\begin{pmatrix}
4&0\\
0&-1
\end{pmatrix}.
\]
Déterminer les valeurs propres et les sous-espaces propres.
\end{exercicebox}

\begin{exercicebox}[2]
Soit :
\[
A=
\begin{pmatrix}
2&1\\
0&2
\end{pmatrix}.
\]
Déterminer les valeurs propres et les sous-espaces propres.  
La matrice est-elle diagonalisable ?
\end{exercicebox}

\begin{exercicebox}[3]
Soit :
\[
A=
\begin{pmatrix}
0&-1\\
1&0
\end{pmatrix}.
\]
Étudier les valeurs propres sur \(\R\), puis sur \(\C\).
\end{exercicebox}

\subsection*{Thème 2 -- Polynôme caractéristique}

\begin{exercicebox}[4]
Calculer le polynôme caractéristique de :
\[
A=
\begin{pmatrix}
1&2&0\\
0&3&0\\
0&0&4
\end{pmatrix}.
\]
Donner les valeurs propres.
\end{exercicebox}

\begin{exercicebox}[5]
Soit :
\[
A=
\begin{pmatrix}
1&1&0\\
0&1&0\\
0&0&2
\end{pmatrix}.
\]
Calculer \(\chi_A\), les sous-espaces propres, puis conclure sur la diagonalisabilité.
\end{exercicebox}

\subsection*{Thème 3 -- Familles de vecteurs propres}

\begin{exercicebox}[6]
Démontrer que trois vecteurs propres associés à trois valeurs propres deux à deux distinctes sont libres.
\end{exercicebox}

\begin{exercicebox}[7]
Soit \(u\in\Lcal(E)\), avec \(\dim(E)=4\).  
On suppose que \(u\) possède quatre valeurs propres distinctes.  
Montrer que \(u\) est diagonalisable.
\end{exercicebox}

\subsection*{Thème 4 -- Multiplicités}

\begin{exercicebox}[8]
Soit :
\[
A=
\begin{pmatrix}
3&1&0\\
0&3&0\\
0&0&5
\end{pmatrix}.
\]
Calculer les multiplicités algébriques et géométriques des valeurs propres.
\end{exercicebox}

\begin{exercicebox}[9]
Donner un exemple de matrice \(2\times2\) dont le polynôme caractéristique est scindé mais qui n'est pas diagonalisable.
\end{exercicebox}

\subsection*{Thème 5 -- Polynômes annulateurs}

\begin{exercicebox}[10]
Soit \(u\in\Lcal(E)\) tel que :
\[
u^2-3u+2\Id_E=0.
\]
Déterminer les valeurs propres possibles de \(u\).  
Que peut-on dire de \(u\) ?
\end{exercicebox}

\begin{exercicebox}[11]
Soit \(p\in\Lcal(E)\) tel que :
\[
p^2=p.
\]
Déterminer les valeurs propres possibles de \(p\), puis montrer que \(p\) est diagonalisable.
\end{exercicebox}

\begin{exercicebox}[12]
Soit \(s\in\Lcal(E)\) tel que :
\[
s^2=\Id_E.
\]
Déterminer les valeurs propres possibles de \(s\), puis montrer que \(s\) est diagonalisable si \(\operatorname{car}(\K)\neq2\).
\end{exercicebox}

%=========================================================
\newpage
\section{Corrigés détaillés des exercices progressifs}

\begin{correctionbox}[1]
La matrice est diagonale.  
Donc :
\[
\chi_A(X)=(X-4)(X+1).
\]
Les valeurs propres sont :
\[
4
\quad\text{et}\quad
-1.
\]

Pour \(\lambda=4\) :
\[
A-4I=
\begin{pmatrix}
0&0\\
0&-5
\end{pmatrix}.
\]
Le système donne :
\[
y=0.
\]
Donc :
\[
E_4=\Vect
\begin{pmatrix}
1\\0
\end{pmatrix}.
\]

Pour \(\lambda=-1\) :
\[
A+I=
\begin{pmatrix}
5&0\\
0&0
\end{pmatrix}.
\]
Le système donne :
\[
x=0.
\]
Donc :
\[
E_{-1}=\Vect
\begin{pmatrix}
0\\1
\end{pmatrix}.
\]
\end{correctionbox}

\begin{correctionbox}[2]
La matrice est triangulaire supérieure.  
Donc :
\[
\chi_A(X)=(X-2)^2.
\]
L'unique valeur propre est :
\[
2.
\]
On calcule :
\[
A-2I=
\begin{pmatrix}
0&1\\
0&0
\end{pmatrix}.
\]
Le système donne :
\[
y=0.
\]
Donc :
\[
E_2=\Vect
\begin{pmatrix}
1\\0
\end{pmatrix}.
\]
Ainsi :
\[
\dim(E_2)=1<2.
\]
La matrice n'est pas diagonalisable.
\end{correctionbox}

\begin{correctionbox}[3]
On calcule :
\[
\chi_A(X)=X^2+1.
\]
Sur \(\R\), il n'y a aucune racine.  
Donc \(A\) n'a pas de valeur propre réelle.

Sur \(\C\) :
\[
X^2+1=(X-i)(X+i).
\]
Les valeurs propres sont :
\[
i
\quad\text{et}\quad
-i.
\]
\end{correctionbox}

\begin{correctionbox}[4]
La matrice est triangulaire supérieure.  
Donc :
\[
\chi_A(X)=(X-1)(X-3)(X-4).
\]
Les valeurs propres sont :
\[
1,\quad3,\quad4.
\]
\end{correctionbox}

\begin{correctionbox}[5]
La matrice est triangulaire supérieure.  
Donc :
\[
\chi_A(X)=(X-1)^2(X-2).
\]

Pour \(\lambda=1\) :
\[
A-I=
\begin{pmatrix}
0&1&0\\
0&0&0\\
0&0&1
\end{pmatrix}.
\]
Le système donne :
\[
y=0,\qquad z=0.
\]
Donc :
\[
E_1=\Vect
\begin{pmatrix}
1\\0\\0
\end{pmatrix}.
\]
Ainsi :
\[
\dim(E_1)=1<2.
\]
La matrice n'est pas diagonalisable.

Pour \(\lambda=2\) :
\[
A-2I=
\begin{pmatrix}
-1&1&0\\
0&-1&0\\
0&0&0
\end{pmatrix}.
\]
Le système donne :
\[
x=0,\qquad y=0.
\]
Donc :
\[
E_2=\Vect
\begin{pmatrix}
0\\0\\1
\end{pmatrix}.
\]
La somme des dimensions vaut :
\[
1+1=2<3.
\]
\end{correctionbox}

\begin{correctionbox}[6]
Soient \(x_1,x_2,x_3\) trois vecteurs propres associés respectivement à des valeurs propres distinctes :
\[
\lambda_1,\lambda_2,\lambda_3.
\]
Supposons :
\[
a_1x_1+a_2x_2+a_3x_3=0.
\]
En appliquant \(u\) :
\[
a_1\lambda_1x_1+a_2\lambda_2x_2+a_3\lambda_3x_3=0.
\]
Multiplions la première relation par \(\lambda_3\) et soustrayons :
\[
a_1(\lambda_1-\lambda_3)x_1+a_2(\lambda_2-\lambda_3)x_2=0.
\]
Les vecteurs \(x_1,x_2\) sont associés à deux valeurs propres distinctes, donc ils sont libres.  
Ainsi :
\[
a_1(\lambda_1-\lambda_3)=0,
\qquad
a_2(\lambda_2-\lambda_3)=0.
\]
Comme les valeurs propres sont distinctes :
\[
a_1=a_2=0.
\]
La relation initiale devient :
\[
a_3x_3=0.
\]
Comme \(x_3\neq0\), on obtient :
\[
a_3=0.
\]
Donc la famille est libre.
\end{correctionbox}

\begin{correctionbox}[7]
On choisit un vecteur propre non nul associé à chacune des quatre valeurs propres distinctes.  
Ces quatre vecteurs sont libres, car ils sont associés à des valeurs propres distinctes.  
Dans un espace de dimension \(4\), une famille libre de \(4\) vecteurs est une base.  
On obtient donc une base de vecteurs propres.  
Ainsi \(u\) est diagonalisable.
\end{correctionbox}

\begin{correctionbox}[8]
La matrice est triangulaire supérieure.  
Donc :
\[
\chi_A(X)=(X-3)^2(X-5).
\]
Les valeurs propres sont :
\[
3
\quad\text{et}\quad
5.
\]

La multiplicité algébrique de \(3\) est :
\[
m_3=2.
\]
La multiplicité algébrique de \(5\) est :
\[
m_5=1.
\]

Pour \(\lambda=3\) :
\[
A-3I=
\begin{pmatrix}
0&1&0\\
0&0&0\\
0&0&2
\end{pmatrix}.
\]
Le système donne :
\[
y=0,
\qquad
z=0.
\]
Donc :
\[
E_3=\Vect
\begin{pmatrix}
1\\0\\0
\end{pmatrix}.
\]
Ainsi la multiplicité géométrique de \(3\) vaut :
\[
\dim(E_3)=1.
\]

Pour \(\lambda=5\) :
\[
A-5I=
\begin{pmatrix}
-2&1&0\\
0&-2&0\\
0&0&0
\end{pmatrix}.
\]
Le système donne :
\[
y=0,
\qquad
x=0.
\]
Donc :
\[
E_5=\Vect
\begin{pmatrix}
0\\0\\1
\end{pmatrix}.
\]
Ainsi :
\[
\dim(E_5)=1.
\]
\end{correctionbox}

\begin{correctionbox}[9]
Un exemple classique est :
\[
A=
\begin{pmatrix}
1&1\\
0&1
\end{pmatrix}.
\]
On a :
\[
\chi_A(X)=(X-1)^2.
\]
Le polynôme caractéristique est scindé.  
Mais :
\[
A-I=
\begin{pmatrix}
0&1\\
0&0
\end{pmatrix},
\]
donc :
\[
E_1=\Vect
\begin{pmatrix}
1\\0
\end{pmatrix}.
\]
La dimension du sous-espace propre est \(1\), strictement inférieure à la multiplicité algébrique \(2\).  
La matrice n'est donc pas diagonalisable.
\end{correctionbox}

\begin{correctionbox}[10]
Le polynôme :
\[
P(X)=X^2-3X+2
\]
annule \(u\).  
On factorise :
\[
P(X)=(X-1)(X-2).
\]
Donc toute valeur propre éventuelle de \(u\) est racine de \(P\), donc appartient à :
\[
\{1,2\}.
\]
Comme \(P\) est scindé à racines simples, \(u\) est diagonalisable.
\end{correctionbox}

\begin{correctionbox}[11]
Si \(p^2=p\), alors :
\[
p^2-p=0.
\]
Donc :
\[
P(X)=X^2-X=X(X-1)
\]
est un polynôme annulateur de \(p\).  
Les valeurs propres possibles sont donc :
\[
0
\quad\text{et}\quad
1.
\]
Le polynôme \(X(X-1)\) est scindé à racines simples, donc \(p\) est diagonalisable.
\end{correctionbox}

\begin{correctionbox}[12]
Si \(s^2=\Id_E\), alors :
\[
s^2-\Id_E=0.
\]
Donc :
\[
P(X)=X^2-1=(X-1)(X+1)
\]
est un polynôme annulateur de \(s\).  
Les valeurs propres possibles sont :
\[
-1
\quad\text{et}\quad
1.
\]
Si \(\operatorname{car}(\K)\neq2\), ces deux racines sont distinctes.  
Le polynôme est alors scindé à racines simples.  
Donc \(s\) est diagonalisable.
\end{correctionbox}

%=========================================================
\newpage
\section{Grand devoir bilan final}

\begin{center}
\Large\textbf{Devoir bilan -- Valeurs propres et vecteurs propres}
\end{center}

\subsection*{Exercice 1 -- Questions de cours}

Soit \(E\) un \(\K\)-espace vectoriel de dimension finie, et soit \(u\in\Lcal(E)\).

\begin{enumerate}
    \item Définir une valeur propre et un vecteur propre de \(u\).
    \item Définir le sous-espace propre \(E_\lambda(u)\).
    \item Montrer que \(E_\lambda(u)\) est un sous-espace vectoriel de \(E\).
    \item Montrer que \(\lambda\) est valeur propre de \(u\) si et seulement si \(E_\lambda(u)\neq\{0\}\).
    \item Montrer que des vecteurs propres associés à des valeurs propres deux à deux distinctes sont libres.
\end{enumerate}

\subsection*{Exercice 2 -- Calcul complet d'éléments propres}

Soit :
\[
A=
\begin{pmatrix}
2&1&0\\
0&2&0\\
0&0&3
\end{pmatrix}.
\]

\begin{enumerate}
    \item Calculer le polynôme caractéristique de \(A\).
    \item Déterminer les valeurs propres de \(A\).
    \item Calculer les sous-espaces propres.
    \item Donner les multiplicités algébriques et géométriques.
    \item La matrice \(A\) est-elle diagonalisable ?
\end{enumerate}

\subsection*{Exercice 3 -- Matrice diagonalisable}

Soit :
\[
B=
\begin{pmatrix}
1&1\\
1&1
\end{pmatrix}.
\]

\begin{enumerate}
    \item Calculer \(\chi_B\).
    \item Déterminer les valeurs propres.
    \item Calculer les sous-espaces propres.
    \item Diagonaliser \(B\).
    \item Calculer \(B^n\) pour \(n\geq1\).
\end{enumerate}

\subsection*{Exercice 4 -- Corps de base}

Soit :
\[
C=
\begin{pmatrix}
0&-1\\
1&0
\end{pmatrix}.
\]

\begin{enumerate}
    \item Calculer \(\chi_C\).
    \item Déterminer le spectre de \(C\) sur \(\R\).
    \item Déterminer le spectre de \(C\) sur \(\C\).
    \item La matrice est-elle diagonalisable sur \(\R\) ? sur \(\C\) ?
\end{enumerate}

\subsection*{Exercice 5 -- Polynôme annulateur}

Soit \(u\in\Lcal(E)\) tel que :
\[
u^3-u=0.
\]

\begin{enumerate}
    \item Donner un polynôme annulateur de \(u\).
    \item Factoriser ce polynôme.
    \item Déterminer les valeurs propres possibles de \(u\).
    \item Montrer que \(u\) est diagonalisable si \(\operatorname{car}(\K)\neq2\).
\end{enumerate}

\subsection*{Exercice 6 -- Projecteur}

Soit \(p\in\Lcal(E)\) tel que :
\[
p^2=p.
\]

\begin{enumerate}
    \item Déterminer les valeurs propres possibles de \(p\).
    \item Montrer que :
    \[
    E_0(p)=\Ker(p).
    \]
    \item Montrer que :
    \[
    E_1(p)=\Imf(p).
    \]
    \item En déduire que :
    \[
    E=\Ker(p)\oplus\Imf(p).
    \]
\end{enumerate}

\newpage

%=========================================================
\section{Correction détaillée du devoir bilan}

\subsection*{Correction de l'exercice 1}

\textbf{1.}  
Un scalaire \(\lambda\in\K\) est une valeur propre de \(u\) s'il existe un vecteur non nul \(x\in E\) tel que :
\[
u(x)=\lambda x.
\]
Un tel vecteur \(x\) est appelé vecteur propre associé à \(\lambda\).

\textbf{2.}  
Le sous-espace propre associé à \(\lambda\) est :
\[
E_\lambda(u)=\Ker(u-\lambda\Id_E).
\]

\textbf{3.}  
Comme :
\[
u-\lambda\Id_E
\]
est une application linéaire, son noyau est un sous-espace vectoriel.  
Donc :
\[
E_\lambda(u)
\]
est un sous-espace vectoriel de \(E\).

\textbf{4.}  
On a :
\[
\lambda\in\Sp(u)
\]
si et seulement s'il existe \(x\neq0\) tel que :
\[
u(x)=\lambda x.
\]
Cela équivaut à :
\[
(u-\lambda\Id_E)(x)=0
\]
avec \(x\neq0\).  
Donc :
\[
E_\lambda(u)\neq\{0\}.
\]

\textbf{5.}  
La démonstration complète est celle du théorème fondamental.  
On part d'une relation linéaire :
\[
\alpha_1x_1+\cdots+\alpha_px_p=0.
\]
On applique \(u\), puis on soustrait \(\lambda_p\) fois la relation initiale :
\[
\alpha_1(\lambda_1-\lambda_p)x_1+\cdots+\alpha_{p-1}(\lambda_{p-1}-\lambda_p)x_{p-1}=0.
\]
Par récurrence, les coefficients sont nuls, puis on conclut que le dernier est nul.  
La famille est libre.

\subsection*{Correction de l'exercice 2}

La matrice :
\[
A=
\begin{pmatrix}
2&1&0\\
0&2&0\\
0&0&3
\end{pmatrix}
\]
est triangulaire supérieure.  
Donc :
\[
\chi_A(X)=(X-2)^2(X-3).
\]
Les valeurs propres sont :
\[
2
\quad\text{et}\quad
3.
\]

Pour \(\lambda=2\) :
\[
A-2I_3=
\begin{pmatrix}
0&1&0\\
0&0&0\\
0&0&1
\end{pmatrix}.
\]
Le système donne :
\[
y=0,
\qquad
z=0.
\]
Donc :
\[
E_2(A)=\Vect
\begin{pmatrix}
1\\0\\0
\end{pmatrix}.
\]
Ainsi :
\[
\dim(E_2)=1.
\]
La multiplicité algébrique de \(2\) est :
\[
m_2=2.
\]

Pour \(\lambda=3\) :
\[
A-3I_3=
\begin{pmatrix}
-1&1&0\\
0&-1&0\\
0&0&0
\end{pmatrix}.
\]
Le système donne :
\[
y=0,
\qquad
x=0.
\]
Donc :
\[
E_3(A)=\Vect
\begin{pmatrix}
0\\0\\1
\end{pmatrix}.
\]
La multiplicité algébrique de \(3\) vaut :
\[
m_3=1,
\]
et :
\[
\dim(E_3)=1.
\]

Comme :
\[
\dim(E_2)+\dim(E_3)=1+1=2<3,
\]
la matrice n'est pas diagonalisable.

\subsection*{Correction de l'exercice 3}

On a :
\[
B=
\begin{pmatrix}
1&1\\
1&1
\end{pmatrix}.
\]
Le polynôme caractéristique est :
\[
\chi_B(X)=
\det
\begin{pmatrix}
X-1&-1\\
-1&X-1
\end{pmatrix}
=
(X-1)^2-1.
\]
Donc :
\[
\chi_B(X)=X^2-2X=X(X-2).
\]
Les valeurs propres sont :
\[
0
\quad\text{et}\quad
2.
\]

Pour \(\lambda=0\), on résout :
\[
BX=0.
\]
On obtient :
\[
x+y=0.
\]
Donc :
\[
E_0=\Vect
\begin{pmatrix}
1\\-1
\end{pmatrix}.
\]

Pour \(\lambda=2\),
\[
B-2I=
\begin{pmatrix}
-1&1\\
1&-1
\end{pmatrix}.
\]
Le système donne :
\[
x=y.
\]
Donc :
\[
E_2=\Vect
\begin{pmatrix}
1\\1
\end{pmatrix}.
\]

On prend :
\[
P=
\begin{pmatrix}
1&1\\
-1&1
\end{pmatrix},
\qquad
D=
\begin{pmatrix}
0&0\\
0&2
\end{pmatrix}.
\]
Alors :
\[
B=PDP^{-1}.
\]
On calcule :
\[
P^{-1}
=
\frac12
\begin{pmatrix}
1&-1\\
1&1
\end{pmatrix}.
\]
Pour \(n\geq1\),
\[
D^n=
\begin{pmatrix}
0&0\\
0&2^n
\end{pmatrix}.
\]
Donc :
\[
B^n=PD^nP^{-1}
=
2^{n-1}
\begin{pmatrix}
1&1\\
1&1
\end{pmatrix}.
\]

\subsection*{Correction de l'exercice 4}

On calcule :
\[
\chi_C(X)=
\det
\begin{pmatrix}
X&1\\
-1&X
\end{pmatrix}
=
X^2+1.
\]

Sur \(\R\), ce polynôme n'a aucune racine.  
Donc :
\[
\Sp_\R(C)=\varnothing.
\]
La matrice n'est donc pas diagonalisable sur \(\R\).

Sur \(\C\), on a :
\[
X^2+1=(X-i)(X+i).
\]
Donc :
\[
\Sp_\C(C)=\{i,-i\}.
\]
Les deux valeurs propres sont distinctes, donc \(C\) est diagonalisable sur \(\C\).

\subsection*{Correction de l'exercice 5}

On a :
\[
u^3-u=0.
\]
Donc :
\[
P(X)=X^3-X
\]
est un polynôme annulateur de \(u\).

On factorise :
\[
P(X)=X(X^2-1)=X(X-1)(X+1).
\]
Les valeurs propres possibles sont donc :
\[
-1,\quad0,\quad1.
\]
Si \(\operatorname{car}(\K)\neq2\), ces racines sont deux à deux distinctes.  
Le polynôme annulateur est donc scindé à racines simples.  
Par le critère des polynômes annulateurs, \(u\) est diagonalisable.

\subsection*{Correction de l'exercice 6}

Comme :
\[
p^2=p,
\]
le polynôme :
\[
P(X)=X^2-X=X(X-1)
\]
annule \(p\).  
Les valeurs propres possibles sont donc :
\[
0
\quad\text{et}\quad
1.
\]

On a :
\[
E_0(p)=\Ker(p-0\Id)=\Ker(p).
\]

Montrons :
\[
E_1(p)=\Imf(p).
\]
Si \(x\in E_1(p)\), alors :
\[
p(x)=x.
\]
Donc \(x\in\Imf(p)\), car \(x=p(x)\).

Réciproquement, si \(x\in\Imf(p)\), alors il existe \(y\in E\) tel que :
\[
x=p(y).
\]
Donc :
\[
p(x)=p(p(y))=p^2(y)=p(y)=x.
\]
Ainsi :
\[
x\in E_1(p).
\]
Donc :
\[
E_1(p)=\Imf(p).
\]

Comme \(p\) est diagonalisable et que ses sous-espaces propres sont \(E_0(p)\) et \(E_1(p)\), on obtient :
\[
E=E_0(p)\oplus E_1(p).
\]
Donc :
\[
E=\Ker(p)\oplus\Imf(p).
\]

%=========================================================
\newpage
\section{Fiche de synthèse finale}

\subsection*{1. Valeur propre et vecteur propre}

\[
\lambda\in\K
\]
est valeur propre de \(u\) si :
\[
\exists x\neq0,\qquad u(x)=\lambda x.
\]

Le vecteur \(x\) est un vecteur propre associé à \(\lambda\).

\subsection*{2. Sous-espace propre}

\[
E_\lambda(u)=\Ker(u-\lambda\Id_E).
\]

\[
\lambda\in\Sp(u)
\Longleftrightarrow
E_\lambda(u)\neq\{0\}.
\]

Les vecteurs propres associés à \(\lambda\) sont :
\[
E_\lambda(u)\setminus\{0\}.
\]

\subsection*{3. Version matricielle}

\[
AX=\lambda X
\Longleftrightarrow
(A-\lambda I_n)X=0.
\]

\[
E_\lambda(A)=\Ker(A-\lambda I_n).
\]

\[
\lambda\in\Sp(A)
\Longleftrightarrow
A-\lambda I_n\text{ non inversible}.
\]

\subsection*{4. Polynôme caractéristique}

\[
\chi_A(X)=\det(XI_n-A).
\]

\[
\lambda\in\Sp(A)
\Longleftrightarrow
\chi_A(\lambda)=0.
\]

Les valeurs propres sont les racines du polynôme caractéristique.

\subsection*{5. Matrices semblables}

Si :
\[
B=P^{-1}AP,
\]
alors :
\[
\chi_B=\chi_A.
\]

Deux matrices semblables ont les mêmes valeurs propres.

\subsection*{6. Vecteurs propres associés à des valeurs propres distinctes}

Des vecteurs propres associés à des valeurs propres deux à deux distinctes sont libres.

Les sous-espaces propres associés à des valeurs propres distinctes sont en somme directe.

\subsection*{7. Multiplicités}

Multiplicité algébrique :
\[
m_\lambda=\text{multiplicité de }\lambda\text{ comme racine de }\chi_u.
\]

Multiplicité géométrique :
\[
\dim(E_\lambda).
\]

Toujours :
\[
1\leq\dim(E_\lambda)\leq m_\lambda.
\]

\subsection*{8. Diagonalisation}

\[
u\text{ diagonalisable}
\Longleftrightarrow
E=\bigoplus_{\lambda\in\Sp(u)}E_\lambda.
\]

Critère dimensionnel :
\[
u\text{ diagonalisable}
\Longleftrightarrow
\sum_{\lambda\in\Sp(u)}\dim(E_\lambda)=\dim(E).
\]

Critère pratique :
\[
u\text{ diagonalisable}
\Longleftrightarrow
\chi_u\text{ est scindé et }\forall\lambda,\ \dim(E_\lambda)=m_\lambda.
\]

Si \(u\) possède \(n\) valeurs propres distinctes en dimension \(n\), alors \(u\) est diagonalisable.

\subsection*{9. Polynômes annulateurs}

Si :
\[
P(u)=0,
\]
alors :
\[
\lambda\in\Sp(u)\Longrightarrow P(\lambda)=0.
\]

Donc :
\[
\Sp(u)\subset\{\text{racines de }P\}.
\]

Si \(P\) est scindé à racines simples et annule \(u\), alors \(u\) est diagonalisable.

\subsection*{10. Cas classiques}

Projecteur :
\[
p^2=p
\Longrightarrow
\Sp(p)\subset\{0,1\}.
\]

Involution :
\[
s^2=\Id
\Longrightarrow
\Sp(s)\subset\{-1,1\}.
\]

Nilpotent :
\[
A^k=0
\Longrightarrow
\Sp(A)\subset\{0\}.
\]

\subsection*{11. Méthode complète}

Pour déterminer les éléments propres d'une matrice \(A\) :
\begin{enumerate}
    \item calculer :
    \[
    \chi_A(X)=\det(XI_n-A);
    \]
    \item résoudre :
    \[
    \chi_A(\lambda)=0;
    \]
    \item pour chaque valeur propre \(\lambda\), résoudre :
    \[
    (A-\lambda I_n)X=0;
    \]
    \item donner une base de chaque sous-espace propre ;
    \item comparer les dimensions aux multiplicités algébriques pour conclure sur la diagonalisabilité.
\end{enumerate}

%=========================================================
\section*{Conclusion pédagogique}
\addcontentsline{toc}{section}{Conclusion pédagogique}

Les valeurs propres et les vecteurs propres permettent de comprendre un endomorphisme à partir des directions qu'il conserve.  
Le passage essentiel est :
\[
u(x)=\lambda x
\quad\Longleftrightarrow\quad
x\in\Ker(u-\lambda\Id_E).
\]

Ainsi, l'étude des valeurs propres transforme un problème géométrique en problème algébrique :
\[
\text{chercher des directions invariantes}
\quad\longleftrightarrow\quad
\text{résoudre des systèmes linéaires}.
\]

Le polynôme caractéristique permet ensuite de trouver les valeurs propres :
\[
\lambda\in\Sp(u)
\Longleftrightarrow
\chi_u(\lambda)=0.
\]

Enfin, l'indépendance des vecteurs propres associés à des valeurs propres distinctes fournit la première idée de la diagonalisation :
\[
\text{assez de vecteurs propres}
\quad\Longrightarrow\quad
\text{base de vecteurs propres}.
\]

Ce chapitre prépare donc naturellement le chapitre complet sur la réduction des endomorphismes.

\end{document}